\documentclass[11pt]{article}
\usepackage{amsmath, amssymb}
\usepackage{geometry}
\usepackage{array}
\usepackage{booktabs}
\geometry{margin=0.7in}
\setlength{\parindent}{0pt}
\setlength{\parskip}{0.35em}
\renewcommand{\arraystretch}{1.15}

\newcommand{\cheatsheettitle}[2]{%
\begin{center}
{\LARGE #1}\\[0.4em]
{\large #2}
\end{center}
}


\begin{document}
\cheatsheettitle{Algebra I Inequalities Cheatsheet}{Module 3}

\section*{Inequality Symbols}
\[
<,\qquad >,\qquad \le,\qquad \ge
\]
\[
\text{open circle: }<\text{ or }>,\qquad
\text{closed circle: }\le\text{ or }\ge
\]

\section*{Solving Rule}
Solve like an equation, except flip the inequality sign when multiplying or dividing by a negative number.
\[
-3x<12 \Rightarrow x>-4
\]

\section*{Interval Notation}
\[
x<5 \Rightarrow (-\infty,5)
\]
\[
x\ge -2 \Rightarrow [-2,\infty)
\]
\[
-1<x\le4 \Rightarrow (-1,4]
\]
Parentheses exclude endpoints. Brackets include endpoints.

\section*{Compound Inequalities}
AND means overlap:
\[
x>2\text{ and }x<7 \Rightarrow 2<x<7
\]
OR means union:
\[
x<-3\text{ or }x\ge5
\]

\section*{Three-Part Inequalities}
\[
-4<2x+6\le10
\]
Subtract $6$ from all three parts:
\[
-10<2x\le4
\]
Divide all three parts by $2$:
\[
-5<x\le2
\]

\section*{Absolute Value Inequalities}
\[
|A|<c \Rightarrow -c<A<c\qquad (c>0)
\]
\[
|A|\le c \Rightarrow -c\le A\le c\qquad (c\ge0)
\]
\[
|A|>c \Rightarrow A<-c\text{ or }A>c\qquad (c\ge0)
\]
\[
|A|\ge c \Rightarrow A\le -c\text{ or }A\ge c\qquad (c\ge0)
\]

\section*{Quick Checks}
\begin{itemize}
  \item Flip the inequality only when multiplying or dividing by a negative.
  \item Use AND for values between two endpoints.
  \item Use OR for values outside two endpoints.
  \item No number has negative absolute value.
\end{itemize}

\end{document}
